% ===== PRÉAMBULE CANONIQUE — à coller VERBATIM en tête de CHAQUE .tex =====
\documentclass[11pt,a4paper]{article}
\usepackage[utf8]{inputenc}\usepackage[T1]{fontenc}\usepackage{lmodern}
\usepackage[french]{babel}
\usepackage{amsmath,amssymb,mathtools}\usepackage{icomma}
\usepackage{siunitx}\sisetup{locale=FR}  % \qty{}{}, \si{}, \ang{}
\usepackage{xcolor}\usepackage{colortbl}
\usepackage{tikz}\usepackage{pgfplots}\pgfplotsset{compat=1.18}
\usepackage[siunitx,europeanresistors]{circuitikz} % circuits (élec.)
\usepackage[most]{tcolorbox}\usepackage{enumitem}\usepackage{array,booktabs}
\usepackage[a4paper,margin=2.2cm]{geometry}
\usetikzlibrary{arrows.meta,calc,angles,quotes,positioning,patterns,%
  backgrounds,shapes.geometric,decorations.pathmorphing,decorations.markings,intersections}
\definecolor{primary}{RGB}{222,49,99}\definecolor{primarydark}{RGB}{150,22,55}
\definecolor{defcol}{RGB}{0,114,178}\definecolor{methcol}{RGB}{52,92,148}
\definecolor{excol}{RGB}{0,135,90}\definecolor{attncol}{RGB}{200,40,55}
\tcbset{boxbase/.style={enhanced,breakable,boxrule=0.9pt,arc=2.5pt,
  left=3mm,right=3mm,top=2.3mm,bottom=2.3mm,fonttitle=\bfseries,coltitle=white}}
% --- boîtes TUTORIEL (noms EXACTS, ne pas renommer) ---
\newtcolorbox{cle}[1][]{boxbase,colback=defcol!6,colframe=defcol,
  title={\textsc{À retenir}\ifx\relax#1\relax\relax\else\textmd{ : \itshape #1}\fi}}
\newtcolorbox{methode}[1][]{boxbase,colback=methcol!7,colframe=methcol,sharp corners,
  title={\textsc{Méthode}\ifx\relax#1\relax\relax\else\textmd{ --- \itshape #1}\fi}}
\newtcolorbox{exemplebox}[1][]{boxbase,colback=excol!5,colframe=excol,
  title={\textsc{Exemple}\ifx\relax#1\relax\relax\else\textmd{ : \itshape #1}\fi}}
\newtcolorbox{piege}[1][]{boxbase,colback=attncol!6,colframe=attncol,
  title={\textsc{Piège}\ifx\relax#1\relax\relax\else\textmd{ : \itshape #1}\fi}}
\newtcolorbox{remarque}[1][]{boxbase,colback=black!4,colframe=black!45,
  title={\textsc{Remarque}\ifx\relax#1\relax\relax\else\textmd{ : \itshape #1}\fi}}
% --- boîtes EXERCICE / CORRIGÉ (noms EXACTS, auto-comptées) ---
\newtcolorbox[auto counter]{exo}[1][]{boxbase,colback=primary!4,colframe=primary,
  title={\textsc{Exercice}~\thetcbcounter\ifx\relax#1\relax\relax\else\textmd{ --- \itshape #1}\fi}}
\newtcolorbox[auto counter]{corrige}[1][]{boxbase,colback=excol!4,colframe=excol,
  title={\textsc{Corrigé}~\thetcbcounter\ifx\relax#1\relax\relax\else\textmd{ --- \itshape #1}\fi}}
\newcommand{\vv}[1]{\vec{#1}}\newcommand{\ds}{\displaystyle}
\newcommand{\R}{\mathbb{R}}\newcommand{\N}{\mathbb{N}}\newcommand{\Z}{\mathbb{Z}}
\begin{document}
\sisetup{per-mode=symbol}

\begin{corrige}[deux positions remarquables]
\begin{enumerate}[label=\alph*)]
  \item Objet à l'infini : $\dfrac{1}{u}\to 0$, donc $\dfrac{1}{v}=\dfrac{1}{f}$ et
        $v=f=+\qty{15}{\cm}$ : image \textbf{réelle} au foyer image $F'$.
  \item $\dfrac{1}{v}=\dfrac{1}{15}-\dfrac{1}{30}=\dfrac{2-1}{30}=\dfrac{1}{30}\Rightarrow v=+\qty{30}{\cm}=2f$ ;
        $\ \gamma=-\dfrac{30}{30}=-1$ : image réelle, renversée, \textbf{même taille}.
\end{enumerate}
\end{corrige}

\begin{corrige}[deux façons de trouver $f$]
\begin{enumerate}[label=\alph*)]
  \item Image réelle : $v=+\qty{1}{\m}$. $\dfrac{1}{f}=\dfrac{1}{1}+\dfrac{1}{1}=2\Rightarrow
        f=\qty{0.5}{\m}=+\qty{50}{\cm}$ (convergente).
  \item Image du même côté que l'objet $\Rightarrow$ \textbf{virtuelle}, $v=-\qty{2.5}{\cm}$.
        $\dfrac{1}{f}=\dfrac{1}{5}+\dfrac{1}{-2{,}5}=0{,}2-0{,}4=-0{,}2\Rightarrow f=-\qty{5}{\cm}$ :
        lentille \textbf{divergente}.
\end{enumerate}
\end{corrige}

\begin{corrige}[l'image agrandie neuf fois]
\begin{enumerate}[label=\alph*)]
  \item Image réelle $\Rightarrow$ renversée $\Rightarrow \gamma=-9$. Or $\gamma=-v/u$, donc
        $v=-\gamma\,u=9\times\qty{1}{\m}=+\qty{9}{\m}$.
  \item $\dfrac{1}{f}=\dfrac{1}{1}+\dfrac{1}{9}=\dfrac{9+1}{9}=\dfrac{10}{9}\Rightarrow
        f=\dfrac{9}{10}=\qty{0.9}{\m}$.
\end{enumerate}
\end{corrige}

\begin{corrige}[la loupe en chiffres]
\begin{enumerate}[label=\alph*)]
  \item $\dfrac{1}{v}=\dfrac{1}{10}-\dfrac{1}{6}=\dfrac{3-5}{30}=-\dfrac{1}{15}\Rightarrow v=-\qty{15}{\cm}$ :
        $v<0$, image \textbf{virtuelle} (\qty{15}{\cm} devant la loupe).
  \item $\gamma=-\dfrac{-15}{6}=+2{,}5$ ; $\ i=\gamma\,o=2{,}5\times\qty{2}{\mm}=+\qty{5}{\mm}$.
        Image \textbf{virtuelle, droite, agrandie} $2{,}5\times$ : c'est le grossissement de la loupe.
\end{enumerate}
\end{corrige}

\begin{corrige}[une lentille divergente]
\begin{enumerate}[label=\alph*)]
  \item $\dfrac{1}{v}=\dfrac{1}{-30}-\dfrac{1}{45}=\dfrac{-3-2}{90}=-\dfrac{5}{90}=-\dfrac{1}{18}
        \Rightarrow v=-\qty{18}{\cm}$ : $v<0$, image \textbf{virtuelle}.
  \item $\gamma=-\dfrac{-18}{45}=+0{,}4$ ; $\ i=0{,}4\times 4=+\qty{1.6}{\cm}$. Image
        \textbf{virtuelle, droite, réduite} : le comportement systématique d'une divergente.
\end{enumerate}
\end{corrige}

\end{document}
